\bookchapter{פרק 54: פרויקט 1 – ספר החשבונות שלי (גרסה מקומית)}

\chapterintro{ברוכים הבאים ל“פרויקט הרציני” הראשון שלכם. ספר החשבונות הזה לא מתחבר לרשת, אבל יש בו את כל האיברים: רישום עסקה, צפייה בפירוט, סיכום חודשי, מחיקה ועריכה. אחרי שתסיימו אותו תעברו מסע שלם של “נתונים לוגיקה מסך”. אל תמהרו להעתיק קוד – קודם חושבים ואז פועלים; זהו העיקרון הקבוע שלנו.}

\booksection{א' חשיבה על דרישות: קודם כול חמש שאלות}

לפני שמתחילים, סדרו את הרעיון עם “חמשת מרכיבי הדרישה”:

\begin{center}
\begin{tabularx}{\linewidth}{|X|X|}
\hline
\textbf{מרכיב} & \textbf{תשובה} \tabularnewline
\hline
Who (למי מיועד) & לשימוש עצמי, וגם בני משפחה וחברים \tabularnewline
What (מה עושים) & רישום הכנסות והוצאות, צפייה בפירוט, סיכום חודשי, סטטיסטיקה לפי קטגוריה \tabularnewline
Why (למה) & תרגול + להבין לאן הכסף הולך \tabularnewline
How (איך) & ‏Kotlin + Compose + Room + MVVM, מקומי לחלוטין בלי שרת עורפי \tabularnewline
Constraint (אילוצים) & להשלים תוך שבוע; תמיכה מינימלית ב-Android 8; גם אם לא יועלה לחנות – לכתוב לפי סטנדרט של אפליקציה בחנות \tabularnewline
\hline
\end{tabularx}
\end{center}

\begin{center}\small\color{gray}חמשת מרכיבי הדרישה של ספר החשבונות\end{center}

ואז פרקו לצעדים קטנים שכל אחד מהם שימושי: (1) בניית שכבת נתונים ב-Room (שיהיה אפשר לשמור רשומה); (2) מסך רשימה (שיהיה אפשר לצפות בפירוט); (3) מסך הזנה (שיהיה אפשר להוסיף רשומה); (4) מחיקה ועריכה; (5) כרטיס סיכום חודשי. אחרי כל צעד, האפליקציה היא “דבר שאפשר ממש להשתמש בו”.

\booksection{ב' שכבת הנתונים: לתת ל“רשומות” בית}

קודם כול מגדירים איך נראית רשומה אחת – השדות צריכים להספיק, אבל בלי חמדנות:

\begin{lstlisting}[language=Kotlin]
@Entity(tableName = "records")
data class Record(
    @PrimaryKey(autoGenerate = true) val id: Long = 0,
    val type: String,        // "income" הכנסה | "expense" הוצאה
    val amount: Double,      // סכום (יואן)
    val category: String,    // קטגוריה: אוכל, תחבורה, משכורת…
    val note: String = "",   // הערה (רשות)
    val date: String         // "2026-08-15"
)
\end{lstlisting}

הוסיפו ב-Dao שתי שאילתות “לצורכי סטטיסטיקה” – זו הנשמה של ספר החשבונות:

\begin{lstlisting}[language=Kotlin]
@Dao
interface RecordDao {
    @Query("SELECT * FROM records ORDER BY date DESC, id DESC")
    fun getAll(): Flow<List<Record>>

    @Insert
    suspend fun insert(record: Record)

    @Delete
    suspend fun delete(record: Record)

    @Query("SELECT COALESCE(SUM(amount),0) FROM records WHERE type='income' AND date LIKE :month || '%'")
    suspend fun monthIncome(month: String): Double   // month בפורמט "2026-08"

    @Query("SELECT COALESCE(SUM(amount),0) FROM records WHERE type='expense' AND date LIKE :month || '%'")
    suspend fun monthExpense(month: String): Double
}
\end{lstlisting}

\begin{tcbinfo}{טיפ: למה הסכום הוא Double?}
אפליקציות פיננסיות קפדניות משתמשות ב-BigDecimal או ב-Int “שסופר באגורות” כדי להימנע משגיאות נקודה צפה (הבעיה הישנה של 0.1+0.2$\neq$0.3). לפרויקט תרגול Double מספיק; אם יום אחד תבנו מוצר פיננסי אמיתי, זכרו לעבור לספירה באגורות – זהו מוקש שדרכו בו אינספור מפתחים לפניכם.

\end{tcbinfo}

\booksection{ג' MVVM: “פירוד” בין המסך ללוגיקה}

לתת למסך לגשת ישירות למסד הנתונים? זה ירוץ, אבל קוד המסך יהפוך לבלגן ארוך ומבאס. ארכיטקטורת MVVM נותנת לכל אחד תפקיד: ‏View (המסך) רק מצייר, ‏ViewModel (מרכז הלוגיקה) מנהל נתונים, ‏Model (Room) מנהל אחסון.

תצוגה (ממשק Compose)

ViewModel (מצב ולוגיקה)

Model (שכבת נתונים Room)

\begin{center}\small\color{gray}איור 5: מבנה MVVM בשלוש שכבות – הנתונים זורמים בכיוון אחד, וכל שכבה נשארת בגבולותיה\end{center}

\begin{lstlisting}[language=Kotlin]
class RecordViewModel(private val dao: RecordDao) : ViewModel() {

    // הופך את ה-Flow למצב שהממשק יכול לצפות בו
    val records = dao.getAll()
        .stateIn(viewModelScope, SharingStarted.WhileSubscribed(), emptyList())

    // פעולות משתמש: קריאה לפונקציות suspend
    fun addRecord(type: String, amount: Double, category: String, note: String, date: String) {
        viewModelScope.launch {
            dao.insert(Record(type = type, amount = amount,
                category = category, note = note, date = date))
        }
    }

    fun deleteRecord(record: Record) {
        viewModelScope.launch { dao.delete(record) }
    }
}
\end{lstlisting}

במסך (Compose) משתמשים כך: val vm: RecordViewModel = viewModel(), ואז val list by vm.records.collectAsState() – ברגע שהנתונים משתנים הרשימה מתרעננת אוטומטית, בלי שתעשו דבר. זו העוצמה של השילוב בין MVVM + Flow + Compose.

\booksection{ד' המסכים: מסך רשימה + מסך הזנה}

במסך הרשימה משתמשים ב-LazyColumn (רשימה בטעינה עצלה) – גם עם הרבה מאוד נתונים, בלי קפאונים:

\begin{lstlisting}[language=Kotlin]
@Composable
fun RecordListScreen(vm: RecordViewModel) {
    val records by vm.records.collectAsState()

    Column {
        // כרטיס הסיכום העליון (הכנסות - הוצאות = יתרה)
        SummaryCard(records)

        // רשימת הרשומות
        LazyColumn {
            items(records, key = { it.id }) { record ->
                RecordRow(record) { vm.deleteRecord(record) }
            }
        }

        // כפתור צף "רישום הוצאה"  פותח דיאלוג הזנה
        FloatingActionButton(onClick = { showAddDialog = true }) {
            Text("+")
        }
    }
}

@Composable
fun RecordRow(record: Record, onDelete: () -> Unit) {
    Row(Modifier.fillMaxWidth().padding(12.dp)) {
        Text(record.category, fontWeight = FontWeight.Bold)
        Spacer(Modifier.weight(1f))
        Text(if (record.type == "income") "+${record.amount}" else "-${record.amount}",
             color = if (record.type == "income") Color(0xFF27AE60) else Color(0xFFC0392B))
    }
}
\end{lstlisting}

להזנה הדרך הפשוטה ביותר היא AlertDialog (תיבת דו-שיח): כמה שדות קלט + מתג הכנסה/הוצאה + תפריט קטגוריות נפתח, ובלחיצה על אישור קוראים ל-vm.addRecord(...).

\begin{tcbwarn}{שימו לב: אל תתעצלו על בדיקת הקלט}
חובה לבדוק את שדה הסכום: ערך ריק, מספר שלילי וקלט לא חוקי כמו “abc” חייבים להיתפס (בדקו עם toDoubleOrNull()). אחת הסיבות הנפוצות ביותר לקריסת אפליקציות היא “להתייחס לקלט של המשתמש כאל ילד טוב”. זכרו: כל קלט של משתמש בודקים כאילו היה אויב.

\end{tcbwarn}

\booksection{ה' עריכה ומחיקה: לאפשר לנתונים “להיכנס וגם לצאת”}

ספר חשבונות שיודע רק להוסיף ולא למחוק מתיישן אחרי שלושה ימים. הוסיפו שתי פעולות:

\begin{lstlisting}[language=Kotlin]
// מוסיפים שתי פונקציות ל-Dao:
@Query("UPDATE records SET amount=:amount, category=:category, note=:note WHERE id=:id")
suspend fun update(id: Long, amount: Double, category: String, note: String)

@Query("DELETE FROM records WHERE id=:id")
suspend fun deleteById(id: Long)

// בממשק: לחיצה ארוכה על פריט פותחת תפריט "עריכה/מחיקה" (ב-Compose: DropdownMenu או AlertDialog)
// לפני מחיקה עדיף דיאלוג אישור — "מחיקה בטעות" היא הדבר שמשתמשי אפליקציות הוצאות שונאים יותר מכול
\end{lstlisting}

שני כללי ניסיון לפרטי האינטראקציה: למחיקה חובה אישור כפול (למנוע החלקת אצבע), ובעריכה כדאי למלא מראש את הערך המקורי (לתקן מילה אחת קל יותר מאשר להקליד הכול מחדש). אל תזלזלו בשני הפרטים האלה – תחושת המוצר בנויה בדיוק מפרטים כאלה.

\booksection{ו' קבלת הפרויקט: בדיקה בשלוש שכבות + שני אתגרים למתקדמים}

\booksection{ה' קבלת הפרויקט: בדיקה בשלוש שכבות + שני אתגרים למתקדמים}

אחרי שהפרויקט מוכן, קבלו אותו לפי שלוש שכבות “תפקוד קצוות ארכיטקטורה”: האם אפשר לרשום, למחוק, והאם הנתונים נשארים אחרי פתיחה מחדש (תפקוד); האם הזנת 0 או מספר שלילי מפילה את האפליקציה, האם סיבוב מסך מאבד נתונים (קצוות); האם המסך והנתונים מעורבבים, האם יש ערכים מקובעים בקוד (ארכיטקטורה).

אחרי שעברתם את קבלת הגרסה הבסיסית, תנו לעצמכם שתי “מנות מתקדמות”:

\begin{enumerate}
  \item סינון חודשי: מוסיפים בראש המסך החלפת חודשים (למשל 2026-08 v), והרשימה מציגה רק את הפירוט של החודש הנוכחי;
  \item תרשים עוגה לפי קטגוריות: משתמשים ב-Canvas של Compose או מוסיפים ספריית תרשימים קלת-משקל, ומציירים את "לאן הכסף הולך" כתרשים עוגה – אחרי זה ספר החשבונות כבר נראה קצת כמו "מוצר אמיתי".
\end{enumerate}

\begin{tcbpractice}{הגיע הזמן להתנסות: השלימו אותו השבוע}
פעלו בסדר “שכבת נתונים ViewModel מסך רשימה מסך הזנה כרטיס סיכום”, צעד אחד בכל יום. נתקעתם? שלחו את הודעת השגיאה כמו שהיא למנוע חיפוש או ל-AI (זוכרים את החלק השביעי, “AI לקידוד בפועל”? בתיאור בעיה אמרו בבהירות “התנהגות צפויה, התנהגות בפועל, שלבי שחזור”). אחרי ההשלמה נהלו בו חשבונות במשך שבוע – כשמשתמשים באפליקציה שבניתם בעצמכם, תחושת ההישג דוחפת אתכם קדימה.

\end{tcbpractice}

\begin{tcbsummary}{סיכום הפרק}
\begin{itemize}
  \item חמשת מרכיבי הדרישה פירוק לצעדים קטנים "שכל אחד שימושי", ורק אז לפעול
  \item שכבת נתונים: ישות Record + Dao (כולל SQL לסטטיסטיקה חודשית) + Database
  \item חלוקת תפקידים ב-MVVM: ‏View רק מצייר, ‏ViewModel מנהל לוגיקה, ‏Model מנהל אחסון
  \item ‏Flow + stateIn + collectAsState: הנתונים משתנים – המסך מתרענן אוטומטית
  \item ‏LazyColumn לרשימה בטעינה עצלה; ‏AlertDialog להזנה
  \item קבלה בשלוש שכבות (תפקוד/קצוות/ארכיטקטורה) + אתגרים למתקדמים (סינון חודשי, תרשים עוגה לפי קטגוריות)
\end{itemize}

\end{tcbsummary}
